% --- Archon physics formalization source begin ---
% archon:physics
% archon:covers PhyXMiniProblems/problem_phyx_mini_0555.lean
% archon:source-report reports/phyx_mini/problem_phyx_mini_0555.source.json

\chapter{Physics problem phyx\_mini\_0555}
\label{ch:PhyXMiniProblems_problem_phyx_mini_0555}

\paragraph{Problem source.}
\#\# Physical scenario

The discovery of the antiproton \textbackslash{}( \textbackslash{}overline\{\textbackslash{}text\{p\}\} \textbackslash{}) (a particle with the same rest energy as a proton, \textbackslash{}( 938 \textbackslash{}text\{ MeV\} \textbackslash{}), but with the opposite electric charge) took place in 1956 through the following reaction:
\textbackslash{}[ \textbackslash{}text\{p\} + \textbackslash{}text\{p\} \textbackslash{}rightarrow \textbackslash{}text\{p\} + \textbackslash{}text\{p\} + \textbackslash{}text\{p\} + \textbackslash{}overline\{\textbackslash{}text\{p\}\} \textbackslash{}]
in which accelerated protons were incident on a target of protons at rest in the laboratory. The minimum incident kinetic energy needed to produce the reaction is called the \textbackslash{}textit\{threshold\} kinetic energy, for which the final particles move together as if they were a single unit as shown in figure.

\#\# Question

Find the threshold kinetic energy to produce antiprotons in this reaction.

\#\# Answer choices

- A: A: \textbackslash{}( 3.901 \textbackslash{}text\{ GeV\} \textbackslash{})

- B: B: \textbackslash{}( 7.581 \textbackslash{}text\{ GeV\} \textbackslash{})

- C: C: \textbackslash{}( 4.364 \textbackslash{}text\{ GeV\} \textbackslash{})

- D: D: \textbackslash{}( 5.628 \textbackslash{}text\{ GeV\} \textbackslash{})

\#\# Figure caption (auxiliary; use the image as primary evidence)

The figure consists of two side-by-side panels separated by a vertical line.

**Left Panel:**
- A 2D coordinate system is displayed with x and y axes labeled.
- There is a single light blue sphere labeled "p" at the origin.
- Another "p" is positioned along the positive x-axis.
- A blue arrow labeled "v" points towards the origin along the negative x-axis, indicating movement.

**Right Panel:**
- Another 2D coordinate system is shown with x and y axes labeled.
- There are three light blue spheres clustered at the origin.
- A single slightly darker sphere labeled "\textbackslash{}(\textbackslash{}bar\{p\}\textbackslash{})" is along the positive x-axis.
- A blue arrow pointing towards the positive x-axis indicates movement.

These panels seem to depict a physical process involving particles labeled as "p" and "\textbackslash{}(\textbackslash{}bar\{p\}\textbackslash{})" with associated movement directions.

\#\# Dataset metadata

Particle Physics; Physical Model Grounding Reasoning, Multi-Formula Reasoning

\paragraph{Recorded answer/context.}
D: D: \textbackslash{}( 5.628 \textbackslash{}text\{ GeV\} \textbackslash{})

\paragraph{Figure/image path.}
/root/proposal\_for\_physic/science-mango/phyx\_mini\_run/phyx\_data/test\_image/555.png

\paragraph{Formalization target.}
create a compiling Lean file with sorry bodies at `PhyXMiniProblems/problem\_phyx\_mini\_0555.lean`. The Lean declarations must preserve the physical quantities, dimensions or dimensional roles, figure labels, governing-law hypotheses, and final relation expressed by this problem.
Use Mathlib/Physlib names found through LeanExplore where available. If a physics API is missing, introduce faithful local abstractions rather than scalar placeholder aliases.

\begin{theorem}[Physics formalization target]
\label{thm:physics:phyx_mini_0555:target}
\lean{PhyXMiniProblems.ProblemPhyXMini0555.problem_phyx_mini_0555}
\uses{def:physics:phyx-mini-0555:phyxminiproblems-problemphyxmini0555-fixedtargetantiprotonsetup, def:physics:phyx-mini-0555:phyxminiproblems-problemphyxmini0555-incidentthresholdkineticenergygev, def:physics:phyx-mini-0555:phyxminiproblems-problemphyxmini0555-matchesfixedtargetthresholdscenario, def:physics:phyx-mini-0555:phyxminiproblems-problemphyxmini0555-matchessuppliedantiprotonfigure, def:physics:phyx-mini-0555:phyxminiproblems-problemphyxmini0555-matchesproblemrestenergydata, def:physics:phyx-mini-0555:phyxminiproblems-problemphyxmini0555-hasphysicalantiprotonproductionparameters, def:physics:phyx-mini-0555:phyxminiproblems-problemphyxmini0555-satisfiesrelativisticthresholdlaws, def:physics:phyx-mini-0555:phyxminiproblems-problemphyxmini0555-matchesdisplayedthresholdchoice}
The assigned autoformalize agent should translate this physics problem into Lean declarations in the covered file, with theorem and lemma proof bodies written as `by sorry`.
\end{theorem}
\begin{proof}
This is an autoformalization task, not a proof task. Produce statements that can later be proved by the physics prover without weakening the physical model.
\end{proof}

% --- Archon declaration topology begin ---
\section{Declaration topology}
\begin{definition}[energy In Electron Volts]
  \label{def:physics:phyx-mini-0555:phyxminiproblems-problemphyxmini0555-energyinelectronvolts}
  \lean{PhyXMiniProblems.ProblemPhyXMini0555.energyInElectronVolts}
  Read a physical energy in electron-volts
\end{definition}
\begin{proof}
  This is the declared model definition of energy In Electron Volts.
\end{proof}
\begin{definition}[energy In Mega Electron Volts]
  \label{def:physics:phyx-mini-0555:phyxminiproblems-problemphyxmini0555-energyinmegaelectronvolts}
  \lean{PhyXMiniProblems.ProblemPhyXMini0555.energyInMegaElectronVolts}
  \uses{def:physics:phyx-mini-0555:phyxminiproblems-problemphyxmini0555-energyinelectronvolts}
  Read a physical energy in mega-electron-volts
\end{definition}
\begin{proof}
  This is the declared model definition of energy In Mega Electron Volts.
\end{proof}
\begin{definition}[energy In Giga Electron Volts]
  \label{def:physics:phyx-mini-0555:phyxminiproblems-problemphyxmini0555-energyingigaelectronvolts}
  \lean{PhyXMiniProblems.ProblemPhyXMini0555.energyInGigaElectronVolts}
  \uses{def:physics:phyx-mini-0555:phyxminiproblems-problemphyxmini0555-energyinelectronvolts}
  Read a physical energy in giga-electron-volts
\end{definition}
\begin{proof}
  This is the declared model definition of energy In Giga Electron Volts.
\end{proof}
\begin{definition}[Particle Species]
  \label{def:physics:phyx-mini-0555:phyxminiproblems-problemphyxmini0555-particlespecies}
  \lean{PhyXMiniProblems.ProblemPhyXMini0555.ParticleSpecies}
  The two particle species occurring in the reaction
\end{definition}
\begin{proof}
  This is the declared model definition of Particle Species.
\end{proof}
\begin{definition}[Electric Charge Sign]
  \label{def:physics:phyx-mini-0555:phyxminiproblems-problemphyxmini0555-electricchargesign}
  \lean{PhyXMiniProblems.ProblemPhyXMini0555.ElectricChargeSign}
  Coarse charge sign, sufficient to record the stated opposite charges
\end{definition}
\begin{proof}
  This is the declared model definition of Electric Charge Sign.
\end{proof}
\begin{definition}[charge Sign]
  \label{def:physics:phyx-mini-0555:phyxminiproblems-problemphyxmini0555-particlespecies-chargesign}
  \lean{PhyXMiniProblems.ProblemPhyXMini0555.ParticleSpecies.chargeSign}
  \uses{def:physics:phyx-mini-0555:phyxminiproblems-problemphyxmini0555-particlespecies, def:physics:phyx-mini-0555:phyxminiproblems-problemphyxmini0555-electricchargesign}
  The proton and antiproton have opposite electric-charge signs
\end{definition}
\begin{proof}
  This is the declared model definition of charge Sign.
\end{proof}
\begin{definition}[Reaction Particle]
  \label{def:physics:phyx-mini-0555:phyxminiproblems-problemphyxmini0555-reactionparticle}
  \lean{PhyXMiniProblems.ProblemPhyXMini0555.ReactionParticle}
  The two incoming and four outgoing particle roles in the reaction
\end{definition}
\begin{proof}
  This is the declared model definition of Reaction Particle.
\end{proof}
\begin{definition}[species]
  \label{def:physics:phyx-mini-0555:phyxminiproblems-problemphyxmini0555-reactionparticle-species}
  \lean{PhyXMiniProblems.ProblemPhyXMini0555.ReactionParticle.species}
  \uses{def:physics:phyx-mini-0555:phyxminiproblems-problemphyxmini0555-particlespecies, def:physics:phyx-mini-0555:phyxminiproblems-problemphyxmini0555-reactionparticle}
  Species of each individually tracked reaction particle
\end{definition}
\begin{proof}
  This is the declared model definition of species.
\end{proof}
\begin{definition}[Is Outgoing]
  \label{def:physics:phyx-mini-0555:phyxminiproblems-problemphyxmini0555-reactionparticle-isoutgoing}
  \lean{PhyXMiniProblems.ProblemPhyXMini0555.ReactionParticle.IsOutgoing}
  \uses{def:physics:phyx-mini-0555:phyxminiproblems-problemphyxmini0555-reactionparticle}
  Predicate selecting exactly the four particles in the final state
\end{definition}
\begin{proof}
  This is the declared model definition of Is Outgoing.
\end{proof}
\begin{definition}[Figure Panel]
  \label{def:physics:phyx-mini-0555:phyxminiproblems-problemphyxmini0555-figurepanel}
  \lean{PhyXMiniProblems.ProblemPhyXMini0555.FigurePanel}
  Initial and threshold-final panels separated by the vertical rule
\end{definition}
\begin{proof}
  This is the declared model definition of Figure Panel.
\end{proof}
\begin{definition}[Figure Axis]
  \label{def:physics:phyx-mini-0555:phyxminiproblems-problemphyxmini0555-figureaxis}
  \lean{PhyXMiniProblems.ProblemPhyXMini0555.FigureAxis}
  The horizontal and vertical axes drawn in each panel
\end{definition}
\begin{proof}
  This is the declared model definition of Figure Axis.
\end{proof}
\begin{definition}[Horizontal Direction]
  \label{def:physics:phyx-mini-0555:phyxminiproblems-problemphyxmini0555-horizontaldirection}
  \lean{PhyXMiniProblems.ProblemPhyXMini0555.HorizontalDirection}
  Horizontal orientation of a velocity arrow in the diagram
\end{definition}
\begin{proof}
  This is the declared model definition of Horizontal Direction.
\end{proof}
\begin{definition}[Figure Color]
  \label{def:physics:phyx-mini-0555:phyxminiproblems-problemphyxmini0555-figurecolor}
  \lean{PhyXMiniProblems.ProblemPhyXMini0555.FigureColor}
  The two visibly different particle colors in the final cluster
\end{definition}
\begin{proof}
  This is the declared model definition of Figure Color.
\end{proof}
\begin{definition}[Figure Particle Glyph]
  \label{def:physics:phyx-mini-0555:phyxminiproblems-problemphyxmini0555-figureparticleglyph}
  \lean{PhyXMiniProblems.ProblemPhyXMini0555.FigureParticleGlyph}
  One raster glyph for each particle displayed in the two panels
\end{definition}
\begin{proof}
  This is the declared model definition of Figure Particle Glyph.
\end{proof}
\begin{definition}[species]
  \label{def:physics:phyx-mini-0555:phyxminiproblems-problemphyxmini0555-figureparticleglyph-species}
  \lean{PhyXMiniProblems.ProblemPhyXMini0555.FigureParticleGlyph.species}
  \uses{def:physics:phyx-mini-0555:phyxminiproblems-problemphyxmini0555-particlespecies, def:physics:phyx-mini-0555:phyxminiproblems-problemphyxmini0555-figureparticleglyph}
  Species represented by a particle glyph
\end{definition}
\begin{proof}
  This is the declared model definition of species.
\end{proof}
\begin{definition}[panel]
  \label{def:physics:phyx-mini-0555:phyxminiproblems-problemphyxmini0555-figureparticleglyph-panel}
  \lean{PhyXMiniProblems.ProblemPhyXMini0555.FigureParticleGlyph.panel}
  \uses{def:physics:phyx-mini-0555:phyxminiproblems-problemphyxmini0555-figurepanel, def:physics:phyx-mini-0555:phyxminiproblems-problemphyxmini0555-figureparticleglyph}
  Panel in which a particle glyph occurs
\end{definition}
\begin{proof}
  This is the declared model definition of panel.
\end{proof}
\begin{definition}[Antiproton Production Figure]
  \label{def:physics:phyx-mini-0555:phyxminiproblems-problemphyxmini0555-antiprotonproductionfigure}
  \lean{PhyXMiniProblems.ProblemPhyXMini0555.AntiprotonProductionFigure}
  \uses{def:physics:phyx-mini-0555:phyxminiproblems-problemphyxmini0555-figurepanel, def:physics:phyx-mini-0555:phyxminiproblems-problemphyxmini0555-figureaxis, def:physics:phyx-mini-0555:phyxminiproblems-problemphyxmini0555-horizontaldirection, def:physics:phyx-mini-0555:phyxminiproblems-problemphyxmini0555-figurecolor, def:physics:phyx-mini-0555:phyxminiproblems-problemphyxmini0555-figureparticleglyph}
  Raster-level information from the image. The positions are dimensionless diagram coordinates, not physical spacetime positions. They are kept apart from the momenta and energies of the collision setup
\end{definition}
\begin{proof}
  This is the declared model definition of Antiproton Production Figure.
\end{proof}
\begin{definition}[Fixed Target Antiproton Setup]
  \label{def:physics:phyx-mini-0555:phyxminiproblems-problemphyxmini0555-fixedtargetantiprotonsetup}
  \lean{PhyXMiniProblems.ProblemPhyXMini0555.FixedTargetAntiprotonSetup}
  \uses{def:physics:phyx-mini-0555:phyxminiproblems-problemphyxmini0555-reactionparticle, def:physics:phyx-mini-0555:phyxminiproblems-problemphyxmini0555-antiprotonproductionfigure}
  The independent physical quantities. fourMomentumGeV stores laboratory frame natural-unit readouts, while both rest energies and the unknown incident threshold kinetic energy remain dimensionful physical quantities. The common final four-velocity is an independent kinematic quantity. The governing-law premise below, rather than this record, relates it to outgoing momenta and imposes its normalization
\end{definition}
\begin{proof}
  This is the declared model definition of Fixed Target Antiproton Setup.
\end{proof}
\begin{definition}[particle Rest Energy]
  \label{def:physics:phyx-mini-0555:phyxminiproblems-problemphyxmini0555-particlerestenergy}
  \lean{PhyXMiniProblems.ProblemPhyXMini0555.particleRestEnergy}
  \uses{def:physics:phyx-mini-0555:phyxminiproblems-problemphyxmini0555-reactionparticle, def:physics:phyx-mini-0555:phyxminiproblems-problemphyxmini0555-fixedtargetantiprotonsetup}
  Rest energy of an individually tracked particle
\end{definition}
\begin{proof}
  This is the declared model definition of particle Rest Energy.
\end{proof}
\begin{definition}[incident Threshold Kinetic Energy Ge V]
  \label{def:physics:phyx-mini-0555:phyxminiproblems-problemphyxmini0555-incidentthresholdkineticenergygev}
  \lean{PhyXMiniProblems.ProblemPhyXMini0555.incidentThresholdKineticEnergyGeV}
  \uses{def:physics:phyx-mini-0555:phyxminiproblems-problemphyxmini0555-energyingigaelectronvolts, def:physics:phyx-mini-0555:phyxminiproblems-problemphyxmini0555-fixedtargetantiprotonsetup}
  GeV readout of the requested threshold kinetic energy
\end{definition}
\begin{proof}
  This is the declared model definition of incident Threshold Kinetic Energy Ge V.
\end{proof}
\begin{definition}[Matches Fixed Target Threshold Scenario]
  \label{def:physics:phyx-mini-0555:phyxminiproblems-problemphyxmini0555-matchesfixedtargetthresholdscenario}
  \lean{PhyXMiniProblems.ProblemPhyXMini0555.MatchesFixedTargetThresholdScenario}
  \uses{def:physics:phyx-mini-0555:phyxminiproblems-problemphyxmini0555-energyingigaelectronvolts, def:physics:phyx-mini-0555:phyxminiproblems-problemphyxmini0555-fixedtargetantiprotonsetup}
  Laboratory-frame facts stated by the prose and shown by the velocity arrows. The target is at rest, the beam travels along positive x, and the threshold cluster moves together along positive x. No magnitude of the unknown beam energy or speed is fixed here
\end{definition}
\begin{proof}
  This is the declared model definition of Matches Fixed Target Threshold Scenario.
\end{proof}
\begin{definition}[Matches Supplied Antiproton Figure]
  \label{def:physics:phyx-mini-0555:phyxminiproblems-problemphyxmini0555-matchessuppliedantiprotonfigure}
  \lean{PhyXMiniProblems.ProblemPhyXMini0555.MatchesSuppliedAntiprotonFigure}
  \uses{def:physics:phyx-mini-0555:phyxminiproblems-problemphyxmini0555-fixedtargetantiprotonsetup}
  Primary-image evidence. It records the axes, two initial proton glyphs, the three-proton/one-antiproton final cluster, the v label, rightward arrows, and the observed relative positions and colors. None of these fields stores the requested threshold kinetic energy
\end{definition}
\begin{proof}
  This is the declared model definition of Matches Supplied Antiproton Figure.
\end{proof}
\begin{definition}[Matches Problem Rest Energy Data]
  \label{def:physics:phyx-mini-0555:phyxminiproblems-problemphyxmini0555-matchesproblemrestenergydata}
  \lean{PhyXMiniProblems.ProblemPhyXMini0555.MatchesProblemRestEnergyData}
  \uses{def:physics:phyx-mini-0555:phyxminiproblems-problemphyxmini0555-energyinmegaelectronvolts, def:physics:phyx-mini-0555:phyxminiproblems-problemphyxmini0555-fixedtargetantiprotonsetup}
  The numerical rest-energy datum and particle/antiparticle equality from the problem prose. 938 MeV is a rest energy, not the unknown kinetic energy
\end{definition}
\begin{proof}
  This is the declared model definition of Matches Problem Rest Energy Data.
\end{proof}
\begin{definition}[Has Physical Antiproton Production Parameters]
  \label{def:physics:phyx-mini-0555:phyxminiproblems-problemphyxmini0555-hasphysicalantiprotonproductionparameters}
  \lean{PhyXMiniProblems.ProblemPhyXMini0555.HasPhysicalAntiprotonProductionParameters}
  \uses{def:physics:phyx-mini-0555:phyxminiproblems-problemphyxmini0555-energyingigaelectronvolts, def:physics:phyx-mini-0555:phyxminiproblems-problemphyxmini0555-fixedtargetantiprotonsetup, def:physics:phyx-mini-0555:phyxminiproblems-problemphyxmini0555-incidentthresholdkineticenergygev}
  Positivity and future-directedness assumptions selecting the physical branch of the relativistic model. They do not determine a numerical kinetic energy
\end{definition}
\begin{proof}
  This is the declared model definition of Has Physical Antiproton Production Parameters.
\end{proof}
\begin{definition}[Satisfies Relativistic Threshold Laws]
  \label{def:physics:phyx-mini-0555:phyxminiproblems-problemphyxmini0555-satisfiesrelativisticthresholdlaws}
  \lean{PhyXMiniProblems.ProblemPhyXMini0555.SatisfiesRelativisticThresholdLaws}
  \uses{def:physics:phyx-mini-0555:phyxminiproblems-problemphyxmini0555-energyingigaelectronvolts, def:physics:phyx-mini-0555:phyxminiproblems-problemphyxmini0555-fixedtargetantiprotonsetup, def:physics:phyx-mini-0555:phyxminiproblems-problemphyxmini0555-particlerestenergy, def:physics:phyx-mini-0555:phyxminiproblems-problemphyxmini0555-incidentthresholdkineticenergygev}
  Mass-shell normalization, four-momentum conservation, the shared final four-velocity threshold condition, and the general relation E\_beam = K\_beam + m\_p c\textasciicircum{}2 in natural-unit readouts. These are governing relations. In particular, no field gives either K\_beam = 6 m\_p c\textasciicircum{}2 or the numerical answer 5.628 GeV; both remain consequences of the threshold kinematics
\end{definition}
\begin{proof}
  This is the declared model definition of Satisfies Relativistic Threshold Laws.
\end{proof}
\begin{definition}[Answer Choice]
  \label{def:physics:phyx-mini-0555:phyxminiproblems-problemphyxmini0555-answerchoice}
  \lean{PhyXMiniProblems.ProblemPhyXMini0555.AnswerChoice}
  Labels of the four threshold-energy choices printed in the problem
\end{definition}
\begin{proof}
  This is the declared model definition of Answer Choice.
\end{proof}
\begin{definition}[displayed Threshold Energy Ge V]
  \label{def:physics:phyx-mini-0555:phyxminiproblems-problemphyxmini0555-displayedthresholdenergygev}
  \lean{PhyXMiniProblems.ProblemPhyXMini0555.displayedThresholdEnergyGeV}
  \uses{def:physics:phyx-mini-0555:phyxminiproblems-problemphyxmini0555-answerchoice}
  Displayed threshold kinetic energy of each answer choice, in GeV
\end{definition}
\begin{proof}
  This is the declared model definition of displayed Threshold Energy Ge V.
\end{proof}
\begin{definition}[Matches Displayed Threshold Choice]
  \label{def:physics:phyx-mini-0555:phyxminiproblems-problemphyxmini0555-matchesdisplayedthresholdchoice}
  \lean{PhyXMiniProblems.ProblemPhyXMini0555.MatchesDisplayedThresholdChoice}
  \uses{def:physics:phyx-mini-0555:phyxminiproblems-problemphyxmini0555-fixedtargetantiprotonsetup, def:physics:phyx-mini-0555:phyxminiproblems-problemphyxmini0555-incidentthresholdkineticenergygev, def:physics:phyx-mini-0555:phyxminiproblems-problemphyxmini0555-answerchoice, def:physics:phyx-mini-0555:phyxminiproblems-problemphyxmini0555-displayedthresholdenergygev}
  A choice is correct when it equals the physical threshold-energy readout
\end{definition}
\begin{proof}
  This is the declared model definition of Matches Displayed Threshold Choice.
\end{proof}
\begin{lemma}[incident Threshold Kinetic Energy Ge V eq six mul rest Energy]
  \label{lem:physics:phyx-mini-0555:phyxminiproblems-problemphyxmini0555-incidentthresholdkineticenergygev-eq-six-mul-restenergy}
  \lean{PhyXMiniProblems.ProblemPhyXMini0555.incidentThresholdKineticEnergyGeV_eq_six_mul_restEnergy}
  \uses{def:physics:phyx-mini-0555:phyxminiproblems-problemphyxmini0555-energyingigaelectronvolts, def:physics:phyx-mini-0555:phyxminiproblems-problemphyxmini0555-fixedtargetantiprotonsetup, def:physics:phyx-mini-0555:phyxminiproblems-problemphyxmini0555-incidentthresholdkineticenergygev, def:physics:phyx-mini-0555:phyxminiproblems-problemphyxmini0555-matchesfixedtargetthresholdscenario, def:physics:phyx-mini-0555:phyxminiproblems-problemphyxmini0555-matchesproblemrestenergydata, def:physics:phyx-mini-0555:phyxminiproblems-problemphyxmini0555-hasphysicalantiprotonproductionparameters, def:physics:phyx-mini-0555:phyxminiproblems-problemphyxmini0555-satisfiesrelativisticthresholdlaws}
  For a fixed proton target, conservation of the invariant mass gives (p\_beam + p\_target)\textasciicircum{}2 = (4 m\_p)\textasciicircum{}2. Together with p\_target = (m\_p, 0), this yields total beam energy 7 m\_p and hence threshold kinetic energy 6 m\_p. This intermediate relation is a conclusion of the governing laws, not an input field
\end{lemma}
\begin{proof}
  The conclusion follows from the governing relations and figure readouts named in the statement of incident Threshold Kinetic Energy Ge V eq six mul rest Energy.
\end{proof}
% --- Archon declaration topology end ---
% --- Archon physics formalization source end ---
